\section{Background}\label{background}

\subsection{Concepts}\label{theory-concepts}

\subsubsection{Motion Planning}\label{motion-planning}
Motion planning is the computational framework for determining a
sequence of valid configurations that move a robot from a start state to
a goal state while interacting with the environment (e.g., avoiding
obstacles). 
\begin{itemize}
    \item \textbf{Configuration Space (C-Space)}: The set of all
mathematically possible robot generalized coordinates \(\mathcal{C}\) (e.g. robot base pose and joint angles). A configuration \(q \in \mathcal{C}\) describes the position and
orientation. The physically (kinematically) feasible configuration space free of collisions is denoted $\mathcal{C}_{free}$.
    \item \textbf{Problem Formulation}: Finding a continuous path
\(\tau: [0, 1] \rightarrow \mathcal{C}_{free}\) such that
\(\tau(0) = q_{start}\) and \(\tau(1) \in \mathcal{C}_{goal}\). If we define an objective function $J$, the task is to find a trajectory $\tau$ such that:
\[
J(\tau^*) = max\{J(\tau): \tau \in \mathcal{C}_{free}\}
\]
Taking the start and goal states into account, this can also be written as:
\[
\tau^* = \operatorname{argmax} \left\{ J(\tau) \;\middle|\; 
\begin{aligned} 
&\tau(0) = q_{\text{start}}, \tau(1) \in \mathcal{C}_{\text{goal}}, \\ 
&\forall t \in [0,1], \tau(t) \in \mathcal{C}_{\text{free}} 
\end{aligned} 
\right\}
\]
\end{itemize}

\subsubsection{Motion Planning as Probabilistic Inference}
Typically, a trajectory \(\tau\) is optimized to minimize a cost \(J(\tau)\). In the
inference perspective (as discussed in \cite{carvalho2023motion}), this is reframed as sampling from a posterior
distribution. An optimality variable \(\mathcal{O}\) (a binary event
behaving like a likelihood) is introduced where
\(p(\mathcal{O}=1|\tau) \propto \exp(-J(\tau))\). The planning problem
is then formulated as sampling suitable trajectories from the posterior:
\[ p(\tau | \mathcal{O}=1) \propto p(\mathcal{O}=1 | \tau) p(\tau) \]
where the prior \(p(\tau)\) enforces smoothness or kinematics (learned
from data), and the likelihood enforces task constraints (collisions,
goals). Diffusion models naturally fit this framework by learning the
gradient of this log-density (score matching). The score of the
posterior can be decomposed:
\[ \nabla_\tau \log p(\tau | \mathcal{O}=1) = \nabla_\tau \log p(\tau) + \nabla_\tau \log p(\mathcal{O}=1 | \tau) \]
The first term is provided by the unconditional diffusion model
(gradient of the data prior), and the second term acts as a
\emph{guidance} signal (gradient of the cost function).


\subsubsection{Generative Models}\label{generative-models}
\begin{itemize}
\item
  \textbf{Variational Autoencoders (VAEs)}: These models learn a
  low-dimensional latent space representation of the data. While fast,
  ``posterior collapse'' or blurry data generation is often observed.
\item
  \textbf{Diffusion Models}\label{basics-of-diffusion-implementation}: A class of likelihood-based models that
  learn to reverse a gradual noising process. Following the notation
  from the survey \cite{wolf2025diffusion}:

  \begin{itemize}
  \tightlist
  \item
    \textbf{Forward Process}: A Markov chain is defined that
    progressively destroys data structure by adding Gaussian noise
    according to a variance schedule \(\beta_k\):
    \[ q(x_{k+1} | x_k) = \mathcal{N}(x_{k+1}; \sqrt{1 - \beta_k} x_k, \beta_k I) \]
  \item
    \textbf{Reverse Process}: The model attempts to approximate the
    intractable posterior \(q(x_{k-1} | x_k)\) using a learned Gaussian
    transition:
    \[ p_\theta(x_{k-1} | x_k) \approx \mathcal{N}(x_{k-1}; \mu_\theta(x_k, k), \Sigma_\theta(x_k, k)) \]
  \item
    \textbf{Training Objective}: In the DDPM formulation, this reduces
    to training a noise prediction network \(\epsilon_\theta\) to
    predict the noise \(\epsilon\) added to \(x_0\), minimizing the
    simplified MSE loss:
    \[ \mathcal{L} = \mathbb{E}_{k, x_0, \epsilon} \left[ || \epsilon - \epsilon_\theta(x_k, k) ||^2 \right] \]
  \item
    \textbf{Sampling}: By chaining these learned transitions, clean
    samples \(x_0\) can be generated starting from pure noise
    \(x_K \sim \mathcal{N}(0, I)\). Optimized samplers like
    \textbf{DDIM} formulate this as a deterministic ODE, allowing for
    fewer steps.
  \end{itemize}
\item
  \textbf{Flow Matching}: A continuous vector field \(v_t(x)\) is
  learned that pushes a probability density path \(p_t\) from noise to
  data. A regression loss on the vector field is minimized:
  \[ \mathcal{L}_{FM} = \mathbb{E}_{t, x_0, x_1} \left[ || v_t(x_t) - (x_1 - x_0) ||^2 \right] \]
  This often results in straighter trajectories in the probability space
  compared to diffusion, enabling faster inference.
\end{itemize}

\subsection{Related Work}\label{related-work}

\subsubsection{Trajectory Optimization}
Trajectory optimization serves as the foundation for modern motion planning, aiming to find optimal paths that respect kinematic and dynamic constraints.
\begin{enumerate}
    \item \textbf{Gradient-based Trajectory Optimization}: Methods such as CHOMP, TrajOpt, and MPC-based solvers formulate planning as a non-convex optimization problem. As noted in \cite{gu2025humanoidlocomotionmanipulationcurrent}, these approaches typically separate the planning of locomotion (base movement) and manipulation (arm movement). A critical limitation is their reliance on high-quality warm starts; without a good initial guess, these optimizers frequently become trapped in local minima due to the complex, non-convex nature of collision avoidance and kinematic constraints.
    \item \textbf{Sampling-based Trajectory Optimization}: Algorithms like RRT* and PRM offer probabilistic completeness but often struggle with the high dimensionality of humanoid robots. Performing whole-body loco-manipulation in real-time is challenging due to sample inefficiency. To mitigate this, search is often performed on a task-dependent subset of degrees of freedom (DoF), keeping other joints fixed, which limits the robot's versatility.
    \item \textbf{Limitations}: Some methods address these challenges by simplifying the problem formulation, for example by fixing the lower body, restricting motion to quasi-static states, or decomposing control into separate upper- and lower-body subsystems. While these strategies reduce computational and modeling complexity, they also increase inaccuracies, limit the robot's range of achievable behaviors, and constrain adaptability in dynamic or unstructured environments.
\end{enumerate}

\subsubsection{Imitation Learning}
Imitation learning (IL) seeks to learn policies directly from expert demonstrations, bypassing the need for explicit cost function engineering.
\begin{itemize}
    \item \textbf{Paradigms}: Common approaches include Behaviour Cloning (BC), Inverse Reinforcement Learning (IRL), Adversarial Motion Priors (AMP), and Gaussian Mixture Models (GMMs). A particularly successful application in fine-grained bimanual manipulation is presented in \cite{zhao2023learningfinegrainedbimanualmanipulation} (ALOHA), which utilizes simple behavior cloning on high-quality teleoperated data.
    \item \textbf{Data Bottleneck}: A major hurdle for humanoid IL is the scarcity of high-quality training data. Kinematic retargeting of human motion capture data is a standard solution, but often introduces artifacts or physical inconsistencies that require further processing.
\end{itemize}

\subsubsection{Generative Models in Robotics}
Generative modeling has emerged as a powerful tool for synthesis across various domains:
\begin{itemize}
    \item \textbf{Variational Autoencoders (VAEs)}: VAEs have been used to learn latent action spaces. The Action Chunking Transformer (ACT) \cite{zhao2023learningfinegrainedbimanualmanipulation} employs a CVAE-based architecture to predict sequences of actions (chunks), effectively modeling temporal dependencies.
    \item \textbf{Character Animation}: In graphics, models like Motion Diffusion Model (MDM) and Diffuse-Cloc generate realistic human motion from text or sparse signals.
    \item \textbf{Robotics and Computer Vision}: Latent diffusion models are increasingly used in computer vision to generate synthetic training data. In robotics, approaches like Diffuser \cite{janner2022planningdiffusionflexiblebehavior}, Diffusion Policy \cite{chi2023diffusion}, and Motion Planning Diffusion (MPD) \cite{carvalho2023motion} treat planning as a generative process, learning to denoise trajectories from random noise.
\end{itemize}

\subsubsection{Model-Based Diffusion Approaches}
Recent advancements have focused on scaling these models to generalist robots.
\begin{itemize}
    \item \textbf{Generalist Robots}: Projects such as Nvidia GR00T, Pi-zero, and BeyondMimic represent the state-of-the-art in applying large-scale generative models to humanoid control, aiming to create "foundation models" for robot motion that can generalize across tasks and embodiments.
\end{itemize}

\subsubsection{Why Learning-Based Planners?}
To overcome the sensitivity
    to initialization and the computational cost of online search,
    learning-based planners leverage large datasets of expert
    trajectories to learn priors. Earlier methods used Gaussian Mixture
    Models (GMMs) or VAEs, but these often struggle to capture the
    complex, multi-modal distributions of valid motions in
    high-dimensional robot state spaces (e.g., a 7-DoF arm or humanoid).
    \textbf{Diffusion models} have emerged as a superior alternative
    because they:

    \begin{enumerate}
    \def\labelenumi{\arabic{enumi}.}
    \tightlist
    \item
      Effectively model \textbf{multi-modal distributions} (capturing
      distinct valid paths, such as going left vs.~right around an
      obstacle) without mode collapse.
    \item
      Can serve as a \textbf{global trajectory prior} for
      optimization-based planners, providing high-quality
      initializations that avoid local minima.
    \item
      Allow for flexible \textbf{conditional sampling} (e.g.,
      conditioning on obstacles or goals) via gradients of the learned
      score function, effectively performing ``planning as inference''
      rather than just playback.
    \end{enumerate}

\subsubsection{Physics-based Animation and Control}
The synthesis of
dynamically realistic motions that are both reactive to the environment
and responsive to user inputs has been a longstanding objective in
physics-based character animation. These capabilities are increasingly
vital for humanoid robotics. The core challenge lies in reconciling
high-level semantic goals---such as interacting with specific
objects---with the stringent low-level physical constraints of balance,
contact dynamics, and actuation limits.

\subsubsection{Kinematic Priors vs.~End-to-End Learning}
Recent advancements
in kinematic motion diffusion models have demonstrated remarkable
capabilities in motion synthesis, conditioned effectively on partial
states, textual descriptions, or even musical rhythms. A key advantage
of these kinematic models is their distinct modularity; they can be
repurposed for unseen downstream tasks through inference-time techniques
such as classifier guidance or inpainting without the need for
retraining.

However, purely kinematic generation often lacks physical fidelity,
leading to artifacts such as foot sliding, floating, or self-collision.
To address this, a hierarchical strategy is adopted in this work: a
tracking policy is employed to follow the generated kinematic plans.
This contrasts with end-to-end approaches that directly predict control
actions to ensure physical realism. While action-prediction models
inherently satisfy dynamic constraints, a lack of versatility is often
observed. Unlike kinematic diffusion models, adaptation to new tasks or
constraints is difficult without expensive retraining, as learned
representations are tightly coupled to specific dynamic contexts and
lack a varying inference-time interface.

% \item
%   \textbf{Reinforcement Learning (RL)}: Note, most RL methods
%   integrating diffusion are offline.

%   \begin{itemize}
%   \tightlist
%   \item
%     \textbf{Limitations of Standard RL}: Learning long-horizon
%     manipulation directly with end-to-end RL is plagued by sample
%     inefficiency and exploration challenges, especially in sparse-reward
%     settings.
%   \item
%     \textbf{Offline RL with Diffusion}: To address data inefficiency,
%     Offline RL leverages static datasets. However, standard methods
%     suffer from distribution shift when the learned policy deviates from
%     the dataset. Diffusion models are now employed in Offline RL (e.g.,
%     \emph{Diffusion-QL}) to better model the behavior policy's
%     distribution, allowing for more stable policy improvement. It has
%     been noted (Wolf et al.~2025) that utilizing the reward signal
%     during training (unlike decoupling it as in \emph{Diffuser}) can
%     improve sample efficiency and policy alignment.
%   \item
%     \textbf{Skill Composition}: A growing trend (Ajay et al.~2023)
%     combines diffusion-based skill learning (lower-level) with
%     high-level RL policies, enabling long-horizon task solving by
%     chaining diffusion-generated motion primitives.
%   \end{itemize}
\item
  \textbf{Imitation Learning (IL)}:

  \begin{itemize}
  \tightlist
  \item
    \textbf{Behavior Cloning (BC)}: Traditionally, control is treated as
    a supervised learning problem (mapping states to actions). Standard
    BC models (like MLPs or Transformers in \emph{ACT}) suffer from the
    ``covariate shift'' problem: small errors accumulate, driving the
    robot away from the training distribution where the policy is
    undefined.
  \item
    \textbf{Diffusion as Expressive Policies}: Diffusion models address
    the multi-modality inherent in expert demonstrations (e.g., an
    expert might go left or right around an obstacle, whereas a unimodal
    BC policy would average them and crash). As discussed in the survey
    \cite{wolf2025diffusion}, diffusion models effectively learn complex
    action distributions and generate smooth trajectories.
  \item
    \textbf{Trajectory vs.~Action Space}: While some methods predict
    single actions, others (\cite{chi2023diffusion}) predict
    sequences (Receding Horizon Control) or entire trajectories (\cite{carvalho2023motion}). Predicting sequences significantly mitigates the error
    accumulation seen in standard BC.
  \item
    \textbf{Flow Matching}: More recently, Flow Matching has emerged as
    a stable alternative to Diffusion for imitation learning (\cite{rouxel2024flow}), offering faster convergence and straighter probability
    flow paths, as noted in the survey.
  \end{itemize}
\item
  \textbf{Diffusion in Robotics}: Diffusion models in robotics primarily
  focus on \textbf{Trajectory Planning and Control}.

  \begin{itemize}
  \item
    \textbf{Trajectory Planning and Control}:

    \begin{itemize}
    \tightlist
    \item
      \emph{Imitation Learning}: \cite{chi2023diffusion} demonstrated stable manipulation by modeling the action
      distribution as a diffusion process, utilizing \emph{receding
      horizon control} where a sequence of actions is predicted but only
      the first few are executed. \cite{carvalho2023motion} explicitly generates collision-free
      trajectories in configuration space between start and goal states,
      operating as a neural motion planner.
    \item
      \emph{Reinforcement Learning / Offline RL}: \cite{janner2022diffuser} proposed a trajectory-level planning framework that
      generates both states and actions. Constraints and goals are
      integrated via ``inpainting'' during the reverse process or
      classifier guidance. Extensions like
      \cite{ajay2023decision} and \cite{wang2023diffusionql}
      incorporate reward maximization directly into the generation
      process, either by conditional sampling on high returns or by using the diffusion model as a policy within a
      Q-learning framework.
    \end{itemize}

  \item
    \textbf{Robotic Grasp Generation}: Methods apply diffusion to learn distributions of valid SE(3) grasp poses \cite{urain2023se3diff} or latent grasp representations \cite{barad2024graspldm}, enabling the generation of diverse and stable grasps.

  \item
    \textbf{Visual Data Augmentation}: Pretrained image-generation models are leveraged to address data scarcity. Techniques include altering backgrounds or rearranging objects to improve the robustness of downstream policies to visual domain shifts.

  \item
    \textbf{Guidance Mechanisms}:
    \begin{itemize}
    \tightlist
    \item
      \emph{Inpainting}: Constraints (e.g., passing through a waypoint)
      are enforced by hard-replacing parts of the trajectory during
      sampling. Instability is sometimes observed (``off-manifold''
      samples).
    \item
      \emph{Classifier Guidance}: A separate learned cost function is
      used to steer generation.
    \item
      \emph{Classifier-Free Guidance (CFG)}: The model is trained with
      and without conditioning labels, and the difference is checked
      during sampling to steer towards the condition. This provides
      robustness but requires retraining for new constraint types.
    \end{itemize}
  \end{itemize}
\item
  \textbf{Kinematic Motion Diffusion and Guidance}: Recent developments
  in the field frame motion synthesis as a conditional diffusion process
  on the kinematic state space.

  \begin{itemize}
  \tightlist
  \item
    \textbf{Manifold Learning}: The manifold of valid poses is
    approximated by the model (e.g., \cite{tevet2023human} for humans,
    \cite{serifi2024robot} for manipulators). Diverse, naturalistic motions
    (e.g., \cite{kalaria2025dreamcontrol}, \cite{zhao2024dartcontrol}) are generated without being strictly
    constrained by dynamics during generation.
  \item
    \textbf{Inference-Time Guidance}: A critical feature is the ability
    to apply guidance (e.g., \cite{xie2024omnicontrol} for spatial constraints)
    during inference. By modifying the denoising step with a gradient
    term \(\nabla_{x_t} J(x_t)\) or through classifier-free guidance,
    the trajectory is steered towards fulfilling constraints (e.g.,
    ``right hand at position \(p\)'') that were not explicitly optimized
    for during training. This adaptability underscores the reasoning for
    choosing a kinematic diffusion planner. \cite{liao2025beyondmimic}, \cite{tevet2024closdclosingloopsimulation}, and \cite{wu2025uniphys} further
    extend this by coupling kinematic priors with robust tracking
    policies or simulation loops for versatile humanoid control. This offers benefits over a
    rigid end-to-end policy---a single foundation model is provided to
    serve multiple downstream tasks.
  \end{itemize}
\item
  \textbf{Sequence Modeling Paradigms}:
    Diffusion Forcing \cite{chen2024diffusion, wu2025uniphys}: A different approach is used in diffusion forcing for time-series modelling, with per-token noise levels. The stability of sequence modeling is combined with the flexibility of diffusion, allowing for ``variable context'' lengths---future states can be forecasted or missing intermediate states between start and goal can be imputed.
\end{itemize}